\chapter{Követelmények}
\section{Funkcionális követelmények}
\subsection*{Felhasználó}
\begin{itemize}
    \item \textbf{Autentikáció:} A felhasználó tud regisztrálni és bejelentkezni
    \item \textbf{Bejelentkezések tárolása:} A felhasználó bejelentkezéseinek tárolása
    \item \textbf{Jogkörök:} Felhasználó jogkörének kezelése: felhasználó, adminisztrátor, vendég
\end{itemize}
\subsection*{Előrehaladás}
\begin{itemize}
    \item \textbf{Rangrendszer:} Elo\footnotemark{} pontrendszerhez hasonló rangrendszer megvalósítása (többjátékos módhoz)
    \item \textbf{Tapasztalati pontok:} XP\footnotemark{} eltárolása minden felhasználóhoz, aminek értéke minden helyes válaszadással nő
\end{itemize}
\subsection*{Játékmenet}
\begin{itemize}
    \item \textbf{Egyjátékos játékmód:} Egyjátékos gyakorló játékmód, feleletválasztós kérdésekkel
    \item \textbf{Többjátékos játékmód:} Többjátékos rangsorolt játékmód, feleletválasztós kérdésekkel, 1 vs. 1 formátumban
    \item \textbf{Meccskeresés:} Két játékos összesorolása, rang alapján
\end{itemize}
\subsection*{Bejelentések}
\begin{itemize}
    \item \textbf{Kérdés bejelentés:} Lehetőség bejelenteni hibásnak vélt kérdéseket
    \item \textbf{Felhasználó bejelentése:} Lehetőség bejelenteni ellenfelet, amennyiben csalást feltételez a felhasználó
\end{itemize}
\footnotetext{Elo (Élő) pont: Prof. Élő Árpád által kifejlesztett pontszámítási rendszer, amely leginkább a sakkban terjedt el, de napjainkban számos online rangsorolt játékmóddal rendelkező játék használja.}
\footnotetext{XP: Tapasztalati pont – egy olyan játékmechanikai változó, amely minden játszma után rendszerint $n$ értékkel növekszik, és nem csökkenhet.}
\newpage
\section{Nem-funkcionális követelmények}
\subsection*{Teljesítmény}
\begin{itemize}
    \item Az oldalak betöltésének csökkentése, Lazy Loading és Eager Loading alkalmazásával
    \item Lehetőség szerinti legkevesebb harmadik féltől származó NPM csomagok felhasználása
\end{itemize}
\subsection*{Biztonság}
\begin{itemize}
    \item Jelszavak titkosított tárolása
    \item Cookie alapú session autentikáció
    \item Service-to-service kommunikáció titkosított tokennel (WebSocket és API között)
    \item Érvényes SSL/TLS tanúsítvány használata
    \item Minimális jogosultságok Docker konténereken belül
\end{itemize}
\subsection*{Használhatóság}
\begin{itemize}
    \item Modern böngészők támogatása
    \item Reszponzív design
    \item HTTP Polling lehetősége, amennyiben a WebSocket nem támogatott a böngésző által
\end{itemize}
\subsection*{Üzemeltetés}
\begin{itemize}
    \item Docker segítségével könnyen üzembe helyezhető webalkalmazás létrehozása
    \item Fejlesztői és éles környezet támogatása, elkülönített konfigurációval
\end{itemize}
\subsection*{Statisztika, csalás elleni intézkedések}
\begin{itemize}
    \item Backend a legkevesebb adatot küldi a kliensnek
    \item Felhasználó pontosságának rögzítése
    \item Válaszadás gyorsaságának rögzítése
    \item JWT tokenek használata a játékmenet során az adatok integritásának védelmére (kérdés- és meccstokenek)
\end{itemize}
\subsection*{Megbízhatóság}
\begin{itemize}
    \item A játszma állapotának megőrzése felhasználó kilépése esetén
    \item Játszma elhagyása esetén ne szakadjon meg a játékmenet a másik félnek
\end{itemize}

\newpage
\section{Használati módok (felhasználó, adminisztrátor, vendég)}
\subsection*{Vendég}
Nem bejelentkezett felhasználó
\begin{itemize}
    \item Főoldal megtekintése
    \item Bejelentkezés, regisztráció
    \item Játékos profilok megtekintése
\end{itemize}

\subsection*{Felhasználó}
Bejelentkezett felhasználó, normál jogkör
\begin{itemize}
    \item Főoldal megtekintése
    \item Játékos profilok megtekintése
    \item Egyjátékos, többjátékos (rangsorolt) módok
    \item Kérdés, ellenfél bejelentése
\end{itemize}

\subsection*{Adminisztrátor}
Bejelentkezett felhasználó, emelt jogkör
\begin{itemize}
    \item Mindazt, amit a felhasználó
    \item Kérdések, válaszok hozzáadása, törlése, szerkesztése
    \item Kérdés -és játékos bejelentések kezelése, bírálása
    \item Felhasználó kitiltása
\end{itemize}

